\begin{abstract}
Dynamic model merging has emerged as a promising paradigm for serving multiple
task-specific adapters in resource-constrained environments by dynamically ensembling
their activations. Recently, several works have proposed highly complex dynamical,
physical, or metaphorical routing architectures---such as continuous-time chemical
reaction kinetics---claiming that classical parametric linear routers catastrophically
fail in low-data regimes. In this work, we present a rigorous methodological audit
and experimental deconstruction of these SOTA claims. Under a skeptical and systematic
evaluation framework, we reveal that the reported failures of classical routers are not
fundamental representational collapses but are rather artifacts of weak evaluation
protocols, specifically a lack of proper initialization and regularization.

By introducing a properly regularized, maximum-entropy zero-initialized classical
linear router, we expose a \emph{small-sample bottleneck} where learning $768$
parameters from only $64$ samples causes severe overfitting, explaining prior
reports of failure and framing SABLE and ChemMerge as highly effective inductive geometric
priors rather than representational panaceas. However, when provided with adequate
calibration budgets ($N_{\text{cal}} = 4000$), classical parametric routers recover
spectacularly. The unregularized Softmax router achieves $76.22\% \pm 0.78\%$ accuracy,
vastly outperforming stateless SABLE ($73.76\% \pm 0.72\%$) and closely approaching the
stateful ChemMerge performance ceiling ($76.90\% \pm 0.68\%$). We isolate a crucial
bias-variance trade-off where strong regularization ($\lambda=10^{-2}$) is essential
under data scarcity but introduces a performance-limiting constraint bias when data is
abundant, capping accuracy at $74.10\% \pm 0.85\%$.

Finally, we deconstruct the claim that stateful trajectory smoothing improves intermediate
representation quality. By measuring layer-wise semantic quality (cosine similarity to correct
task prototypes), we demonstrate that ChemMerge's "representational lag" is not a defect,
but rather acts as a beneficial temporal low-pass filter (closed-loop stateful inertia).
This control-theoretic feedback loop stabilizes ensembling trajectories under heavy activation
noise, explaining its superior performance ceiling.
\end{abstract}